\section{Arquitetura e Especificações Técnicas}

A concepção do SIRDA (Sistema de Inteligência e Resposta a Desastres Ambientais) afasta-se do uso exclusivo de bancos de dados relacionais clássicos, adotando um ecossistema distribuído de \textit{Big Data} construído inteiramente sobre ferramentas de código aberto (\textit{open-source}). Esta escolha técnica provê escalabilidade ao acesso analítico e garante a capacidade de processar volumes massivos de anomalias térmicas, cruzando essas informações com geometrias complexas de infraestrutura em uma janela de tempo crítica. A arquitetura foi estruturada em torno das seguintes tecnologias:

\begin{itemize}
    \item \textbf{Apache Kafka:} Atua como o sistema nervoso central de mensageria, absorvendo picos severos de dados durante períodos críticos de queimadas, promovendo o desacoplamento entre a extração dos dados e a ingestão no banco analítico.
    \item \textbf{Apache Druid:} O motor OLAP (\textit{Online Analytical Processing}) escolhido pela sua capacidade nativa de ingerir fluxos de dados diretamente do Kafka (\textit{Streaming Ingestion}) e fornecer respostas analíticas em escala de baixa latência (poucos segundos).
    \item \textbf{PostgreSQL:} Empregado como o banco relacional de suporte, atuando estritamente como \textit{Metadata Store} para o Druid e para o sistema de visualização gráfica.
    \item \textbf{Dagster:} Orquestrador central que rege o paradigma DataOps do projeto. Ele garante a execução limpa, tipada e rastreável dos processos de enriquecimento espacial e publicação de mensagens no Kafka.
    \item \textbf{Apache Superset:} Camada visual que materializa o Nível 3 da Consciência Situacional, entregando painéis gerenciais e topológicos aos tomadores de decisão.
\end{itemize}

Para assegurar a reprodutibilidade e a portabilidade do sistema, a infraestrutura foi inteiramente conteinerizada via Docker. A topologia orquestrada no \textit{docker-compose} reflete fielmente um ambiente de produção resiliente. Destaca-se que o banco de dados analítico Apache Druid não foi instanciado como um serviço monolítico, mas sim desmembrado em seus microsserviços fundamentais: \textit{Coordinator}, \textit{Broker}, \textit{Historical}, \textit{MiddleManager} e \textit{Router}. Essa granularidade permite que a gestão de metadados, o roteamento de consultas e o processamento de dados ocorram em instâncias independentes, mitigando gargalos e falhas em cascata.

Além da separação lógica, a implantação local adota restrições físicas de hardware para mimetizar limitações produtivas reais. Foram aplicados limites explícitos de uso de processamento e memória RAM para cada serviço (\textit{deploy resources limits}) na configuração dos contêineres. Paralelamente, os parâmetros da Máquina Virtual Java (JVM) de cada componente do Druid foram estritamente ajustados, por meio das variáveis \textit{JAVA\_OPTS}, definindo precisamente os tamanhos do \textit{heap} (\textit{-Xms} e \textit{-Xmx}) e da memória direta (\textit{MaxDirectMemorySize}). Esses controles evitam o esgotamento prematuro de recursos da máquina hospedeira (\textit{Out of Memory}), assegurando a estabilidade ininterrupta do \textit{pipeline} de resposta a desastres, conforme sumarizado na Tabela \ref{tab:recursos_docker}.

\begin{table}[H]
    \centering
    \caption{Alocação de Recursos e Parâmetros JVM (Docker).}
    \vspace{0.2cm}
    \resizebox{\textwidth}{!}{%
    \begin{tabular}{|l|c|c|c|c|}
        \hline
        \textbf{Serviço} & \textbf{vCPUs} & \textbf{RAM Total} & \textbf{JVM Heap} (\textit{Min/Max}) & \textbf{JVM Direct Mem.} \\
        \hline
        Zookeeper & 0.25 & 256M & 128M / 128M & - \\
        Apache Kafka & 0.75 & 512M & 256M / 256M & - \\
        PostgreSQL & 0.50 & 256M & - & - \\
        Druid Coordinator & 0.50 & 384M & 128M / 256M & - \\
        Druid Broker & 1.0 & 768M & 256M / 512M & 512M \\
        Druid Historical & 1.0 & 768M & 256M / 512M & 512M \\
        Druid MiddleManager & 1.0 & 1024M & 128M / 128M\textsuperscript{*} & 128M\textsuperscript{*} \\
        Druid Router & 0.25 & 384M & 128M / 128M & - \\
        Apache Superset & 1.0 & 768M & - & - \\
        Dagster & 0.75 & 256M & - & - \\
        \hline
    \end{tabular}%
    }
    
    \vspace{0.1cm}
    \raggedright
    \footnotesize{\textsuperscript{*} O \textit{MiddleManager} também aloca instâncias filhas (\textit{Peon Workers}) com Heap de 512M / 512M.}
    
    \fonte{Elaborada pelos autores.}
    \label{tab:recursos_docker}
\end{table}

Os valores estipulados na Tabela \ref{tab:recursos_docker} representam a configuração \textbf{Baseline} estática da arquitetura. Conforme detalhado nos resultados, durante a bateria de estresse de vazão analítica, essas cotas de processamento para os nós \textit{Broker} e \textit{Historical} foram sistematicamente escalonadas até 2.0 vCPUs para mapear os ganhos de eficiência do sistema sob pressão.

A assimetria observada reflete diretamente o perfil analítico de cada nó. O \textit{MiddleManager} aloca a maior cota absoluta de RAM por orquestrar o \textit{streaming} volátil e o agrupamento inicial de dados antes da indexação em disco. Em contrapartida, instâncias focadas em pesquisa complexa (\textit{Broker} e \textit{Historical}) maximizam a alocação de vCPU e Memória Direta, preservando o \textit{Heap} padrão da JVM para evitar pausas prolongadas de \textit{Garbage Collection} durante os cruzamentos espaciais. Os nós de gestão (\textit{Coordinator} e \textit{Router}) operam com capacidade reduzida, pois não realizam operações matemáticas diretas.

\subsection{Modelagem Analítica: Esquema Constelação}

Para suprir os requisitos analíticos de baixa latência do Apache Druid, a modelagem de dados afasta-se da normalização tradicional e adota um \textit{Esquema Constelação} altamente desnormalizado. Neste modelo, as anomalias térmicas brutas e os agrupamentos de desastres formam as duas Tabelas Fato centrais, cercadas por Dimensões de busca contendo as chaves geográficas e infraestruturais pré-computadas, otimizando as agregações interativas na interface visual (Figura \ref{fig:diagrama_er}).

\begin{figure}[H]
    \centering
    \resizebox{\textwidth}{!}{
    \begin{tikzpicture}[
        table/.style={
            rectangle split,
            rectangle split parts=2,
            draw=black!70,
            thick,
            font=\small,
            align=left,
            rounded corners=3pt,
            rectangle split part fill={gray!15, white}
        },
        fato/.style={
            table,
            font=\small\bfseries
        },
        dimensao/.style={
            table,
            font=\footnotesize
        },
        relation/.style={
            -latex,
            very thick,
            draw=black!50
        },
        background_box/.style={
            rectangle, draw=gray!40, thick, dashed, fill=gray!3, rounded corners=6pt, inner sep=15pt
        },
        label_text/.style={font=\scriptsize\sffamily\bfseries, fill=gray!5, draw=gray!40, rounded corners=2pt, inner sep=3pt}
    ]

    % --- FATOS (EM CIMA) ---
    \node[fato] (fato_focos) {
        FATO\_FOCOS\_INCENDIO
        \nodepart{two}
        \normalfont
        PK: foco\_id \\
        \textit{timestamp} \\
        M: satelite \\
        M: latitude, longitude \\
        M: frp (Potência Radiativa) \\
        M: risco\_ons\_calc
    };

    \node[fato, right=3cm of fato_focos] (fato_eventos) {
        FATO\_EVENTOS\_INCENDIO
        \nodepart{two}
        \normalfont
        PK: evento\_id \\
        \textit{timestamp\_inicio} \\
        \textit{timestamp\_fim} \\
        M: densidade \\
        M: risco\_cascata\_global \\
        M: area\_estimada\_km2
    };

    % --- DIMENSOES (EM BAIXO) ---
    % Linha 1 de Dimensões (4 tabelas)
    \node[dimensao, below left=3cm and -1.5cm of fato_focos] (dim_mun) {
        \textbf{DIM\_MUNICIPIOS}
        \nodepart{two}
        PK: codigo\_ibge \\
        estado\_sigla \\
        municipios\_nome
    };

    \node[dimensao, right=0.6cm of dim_mun] (dim_ti) {
        \textbf{DIM\_TERRAS\_INDIGENAS}
        \nodepart{two}
        PK: codigo\_ti \\
        nome\_terra\_indigena \\
        area\_hectares
    };

    \node[dimensao, right=0.6cm of dim_ti] (dim_hosp) {
        \textbf{DIM\_HOSPITAIS}
        \nodepart{two}
        PK: codigo\_cnes \\
        nome\_hospital \\
        tipo\_unidade
    };

    \node[dimensao, right=0.6cm of dim_hosp] (dim_rod) {
        \textbf{DIM\_RODOVIAS}
        \nodepart{two}
        PK: id\_trecho\_rodoviario \\
        codigo\_rodovia \\
        jurisdicao
    };

    % Linha 2 de Dimensões (3 tabelas centralizadas)
    \node[dimensao, below=1.5cm of dim_ti, xshift=-1cm] (dim_lt) {
        \textbf{DIM\_LINHAS\_TRANSMISSAO}
        \nodepart{two}
        PK: id\_linha\_transmissao \\
        nome\_linha \\
        tensao\_kv
    };

    \node[dimensao, right=0.6cm of dim_lt] (dim_sub) {
        \textbf{DIM\_SUBESTACOES}
        \nodepart{two}
        PK: id\_subestacao \\
        nome\_subestacao \\
        capacidade\_mw
    };
    
    \node[dimensao, right=0.6cm of dim_sub] (dim_sis) {
        \textbf{DIM\_SISTEMAS\_ISOLADOS}
        \nodepart{two}
        PK: id\_sistema\_isolado \\
        nome\_sistema\_isolado \\
        prioridade
    };

    % --- GRUPO BACKGROUND (DIMENSÕES COMPARTILHADAS) ---
    \begin{scope}[on background layer]
        \coordinate (bg_top) at ([yshift=0.8cm]dim_mun.north);
        \node[background_box, fit=(bg_top) (dim_mun) (dim_ti) (dim_hosp) (dim_rod) (dim_lt) (dim_sub) (dim_sis)] (bg_dimensoes) {};
        \node[anchor=north west, font=\bfseries\color{black!70}, xshift=5pt, yshift=-5pt] at (bg_dimensoes.north west) {Dimensões Compartilhadas (Esquema Constelação)};
    \end{scope}

    % --- RELACIONAMENTOS (Setas Reduzidas) ---
    % Apenas 2 setas direcionadas ao Container de Dimensões
    \draw[relation] (fato_focos.south) -- node[label_text, right=3pt] {Foreign Keys} (bg_dimensoes.north -| fato_focos.south);
    \draw[relation] (fato_eventos.south) -- node[label_text, left=3pt] {Foreign Keys} (bg_dimensoes.north -| fato_eventos.south);

    \end{tikzpicture}
    }
    \caption{Modelo Entidade-Relacionamento (Constelação).}
    \fonte{Elaborada pelos autores.}
    \label{fig:diagrama_er}
\end{figure}


\subsection{Orquestração Orientada a Dados (Dagster)}

Para orquestrar o cruzamento das malhas elétricas e inferências de ciência de redes, o projeto adota uma abordagem estritamente focada em dados por meio da abstração de \textit{Software-Defined Assets} (SDA) do Apache Dagster. Diferentemente de orquestradores tradicionais focados em fluxos de tarefas genéricas, o modelo implementado estabelece um grafo declarativo de linhagem explícita (\textit{Lineage}). 

Como ilustrado na Figura \ref{fig:diagrama_dagster}, a arquitetura não utiliza agendamentos temporais órfãos. A extração da infraestrutura linear (\texttt{aneel\_ons\_data}) figura como a raiz do grupo \texttt{infra\_ons}. Essa malha base torna-se pré-requisito funcional absoluto (\texttt{deps=[...]}) para a segmentação subsequente das tabelas de Subestações e Linhas de Transmissão. Somente após a materialização bem-sucedida dessas dependências físicas é que o Dagster autoriza a execução dos algoritmos analíticos pesados do grupo \texttt{science\_models}, como o cálculo topológico \texttt{network\_centrality\_model}. 

Isso garante idempotência e tipagem granular: falhas em rotinas de extração bloqueiam instantaneamente a propagação de anomalias para os motores de agregação espacial (NetworkX e Scikit-Learn), blindando a modelagem arquitetural contra estados incompletos.

\begin{figure}[H]
    \centering
    \resizebox{\textwidth}{!}{
    \begin{tikzpicture}[
        font=\sffamily\small,
        %% -- ESTILOS --
        asset/.style={rectangle, draw=black!60, thick, fill=white, rounded corners=4pt, drop shadow={opacity=0.15, shadow xshift=2pt, shadow yshift=-2pt}, align=center, inner sep=6pt, minimum width=3.5cm, minimum height=1.0cm},
        relation/.style={-Stealth, thick, draw=gray!80},
        %% -- GRUPOS --
        bg_infra/.style={rectangle, draw=cyan!60!blue, thick, dashed, fill=cyan!5, rounded corners=6pt, inner sep=12pt},
        bg_science/.style={rectangle, draw=purple!60!blue, thick, dashed, fill=purple!5, rounded corners=6pt, inner sep=12pt},
        bg_dim/.style={rectangle, draw=gray!50, thick, dashed, fill=gray!8, rounded corners=6pt, inner sep=12pt},
        lbl/.style={font=\sffamily\bfseries\scriptsize, text=black, fill=white, draw=gray!40, rounded corners=2pt, inner sep=2pt}
    ]

    % ==========================================
    % INFRA_ONS
    % ==========================================
    \node[asset] at (0, 0) (aneel) {\textbf{aneel\_ons\_data}\\(Python)};
    
    \node[asset] at (6.5, 2.0) (dim_lt) {\textbf{dim\_linhas\_transmissao}\\(Python)};
    \node[asset] at (6.5, 0) (dim_sub) {\textbf{dim\_subestacoes}\\(Python)};
    \node[asset] at (6.5, -2.0) (dim_sis) {\textbf{dim\_sistemas\_isolados}\\(Python)};

    % ==========================================
    % SCIENCE_MODELS
    % ==========================================
    \node[asset] at (13.5, 1.0) (network) {\textbf{network\_centrality\_model}\\(NetworkX)};
    \node[asset] at (13.5, -2.0) (dbscan) {\textbf{st\_dbscan\_fire\_events}\\(Scikit-Learn)};

    % ==========================================
    % DIMENSIONS (Static)
    % ==========================================
    \node[asset] at (0, -5.5) (dim_mun) {\textbf{dim\_municipios}\\(Python)};
    \node[asset] at (4.5, -5.5) (dim_hosp) {\textbf{dim\_hospitais}\\(Python)};
    \node[asset] at (9.0, -5.5) (dim_ti) {\textbf{dim\_terras\_indigenas}\\(Python)};
    \node[asset] at (13.5, -5.5) (dim_rod) {\textbf{dim\_rodovias\_federais}\\(Python)};

    % ==========================================
    % BACKGROUNDS
    % ==========================================
    \begin{scope}[on background layer]
        % Infra
        \node[bg_infra, fit={([shift={(-0.5,1.0)}]aneel.west|-dim_lt.north) ([shift={(0.5,-0.5)}]dim_lt.east|-dim_sis.south)}] (bg_infra_node) {};
        \node[anchor=north west, font=\sffamily\color{black}, xshift=6pt, yshift=-6pt] at (bg_infra_node.north west) {\textbf{Group: infra\_ons}};

        % Science
        \node[bg_science, fit={([shift={(-0.5,1.0)}]network.north west) ([shift={(0.5,-0.5)}]dbscan.south east)}] (bg_science_node) {};
        \node[anchor=north west, font=\sffamily\color{black}, xshift=6pt, yshift=-6pt] at (bg_science_node.north west) {\textbf{Group: science\_models}};

        % Dimensions
        \node[bg_dim, fit={([shift={(-0.5,1.0)}]dim_mun.north west) ([shift={(0.5,-0.5)}]dim_rod.south east)}] (bg_dim_node) {};
        \node[anchor=north west, font=\sffamily\color{black}, xshift=6pt, yshift=-6pt] at (bg_dim_node.north west) {\textbf{Group: dimensions}};
    \end{scope}

    % ==========================================
    % SETAS DE LINHAGEM (DEPENDÊNCIAS)
    % ==========================================
    \draw[relation] (aneel.east) to[out=0, in=180] (dim_lt.west);
    \draw[relation] (aneel.east) to[out=0, in=180] (dim_sub.west);
    \draw[relation] (aneel.east) to[out=0, in=180] (dim_sis.west);

    \draw[relation] (dim_lt.east) to[out=0, in=180] (network.west);
    \draw[relation] (dim_sub.east) to[out=0, in=180] (network.west);

    \end{tikzpicture}
    }
    \caption{Grafo de Dependência Declarativa (Software-Defined Assets).}
    \label{fig:diagrama_dagster}
    \fonte{Elaborada pelos autores.}
\end{figure}




\subsection{Data Flow End-to-End (Streaming e OLAP)}

A Figura \ref{fig:diagrama_ingestao} ilustra a primeira fase do fluxo contínuo de processamento da arquitetura \textit{Medallion}. A extração inicia-se a partir das coordenadas brutas de satélites e das fontes dimensionais. Em vez de persistir o dado bruto inicialmente em disco, a principal premissa arquitetural adotada foi a paralelização em nível de mensageria através do Produtor Kafka, permitindo o enriquecimento em tempo real.

Esse fluxo contínuo (\textit{Streaming Dual-Branch}) bifurca as informações imediatamente: o \textit{Ramo dos Focos} analisa os registros pontuais de forma quase instantânea, enquanto o \textit{Ramo dos Eventos} acumula janelas temporais para a agregação poligonal no \textit{DBSCAN}. Ambas as esteiras invocam o \textit{Motor Espacial} simultaneamente e empacotam seus resultados sob contratos rigidamente tipados (\texttt{FocoFatoPayload} e \texttt{EventoFatoPayload}).

\begin{figure}[H]
    \centering
    \resizebox{\textwidth}{!}{
    \begin{tikzpicture}[
        font=\sffamily\small,
        %% -- ESTILOS DE COMPONENTES --
        process/.style={rectangle, draw=black!60, thick, fill=white, rounded corners=4pt, drop shadow={opacity=0.15, shadow xshift=2pt, shadow yshift=-2pt}, align=center, inner sep=6pt, minimum width=2.8cm, minimum height=1.3cm},
        database/.style={cylinder, cylinder uses custom fill, cylinder body fill=white, cylinder end fill=gray!10, shape border rotate=90, draw=black!60, thick, drop shadow={opacity=0.15, shadow xshift=2pt, shadow yshift=-2pt}, aspect=0.25, align=center, inner sep=5pt, minimum width=2.8cm, minimum height=1.5cm},
        payload/.style={ellipse, draw=black!60, thick, fill=white, drop shadow={opacity=0.15, shadow xshift=2pt, shadow yshift=-2pt}, align=center, inner sep=4pt, minimum width=2.5cm, minimum height=1.1cm},
        source/.style={tape, tape bend top=none, draw=black!60, thick, fill=white, drop shadow={opacity=0.15, shadow xshift=2pt, shadow yshift=-2pt}, align=center, inner sep=6pt, minimum width=2.8cm, minimum height=1.3cm},
        queue/.style={chamfered rectangle, chamfered rectangle corners={north east, south east}, chamfered rectangle xsep=0.4cm, draw=blue!50, thick, fill=cyan!5, drop shadow={opacity=0.15, shadow xshift=2pt, shadow yshift=-2pt}, align=center, inner sep=6pt, minimum width=2.8cm, minimum height=1.3cm},
        %% -- ESTILOS DE BACKGROUND (MEDALLION) --
        bg_bronze/.style={rectangle, draw=orange!40!brown, thick, dashed, fill=orange!5, rounded corners=6pt, inner sep=12pt},
        bg_silver/.style={rectangle, draw=gray!50, thick, dashed, fill=gray!8, rounded corners=6pt, inner sep=12pt},
        bg_gold/.style={rectangle, draw=yellow!60!orange, thick, dashed, fill=yellow!10, rounded corners=6pt, inner sep=12pt},
        %% -- ESTILOS DE ARESTAS --
        relation/.style={-Stealth, thick, draw=gray!70},
        dashed_rel/.style={Stealth-Stealth, thick, dashed, draw=gray!80},
        lbl/.style={font=\sffamily\bfseries\scriptsize, text=black, fill=white, draw=gray!40, rounded corners=2pt, inner sep=3pt}
    ]

    %% ============================================================
    %% GRADE DE LAYOUT OTIMIZADA (Normalizada na Origem X=0)
    %% ============================================================

    % ── COLUNA 1 : RAW / BRONZE ──
    \node[source] at (0, 0)    (api)     {\textbf{API Focos INPE}\\(CSV / API)};
    \node[source] at (0, -4.5) (dim_raw) {\textbf{Fontes Dimensionais}\\(FUNAI, IBGE, CSVs)};

    % ── COLUNA 2 : PRODUTOR KAFKA ──
    \node[process] at (6.0, 0) (prod) {\textbf{Produtor Kafka}\\(Ingestão Python)};

    % ── COLUNA 3 : FORK (Ramos) ──
    \node[process] at (11.0, 4.5) (proc_focos) {\textbf{Processador}\\Espacial};
    \node[queue] at (11.0, -4.5) (buffer) {\textbf{Buffer Temporário}};

    % ── COLUNA 4 : PROCESSAMENTO / MOTOR (Centro) ──
    \node[process] at (17.5, -4.5) (dbscan) {\textbf{Eventizador}\\(DBSCAN)};
    \node[process] at (17.5, 0) (motor) {\textbf{Motor Espacial}\\(Dimensões)};

    % ── COLUNA 5 : PAYLOADS ──
    \node[payload] at (23.0, 4.5)  (pay_focos)   {\textbf{FocoFato}\\Payload};
    \node[payload] at (23.0, -4.5) (pay_eventos) {\textbf{EventoFato}\\Payload};

    % ── COLUNA 6 : TÓPICOS KAFKA (SAÍDA) ──
    \node[queue] at (27.5, 4.5)  (topico_focos)   {\textbf{Tópico Kafka:}\\focos};
    \node[queue] at (27.5, -4.5) (topico_eventos) {\textbf{Tópico Kafka:}\\eventos};

    % ── SILVER LAYER (Faixa inferior isolada) ──
    \node[process]  at (6.0, -11)  (silver_etl)  {\textbf{ETL Dimensões}\\(Shapely/GeoPandas)};
    \node[database] at (14.0, -11) (silver_data) {\textbf{Dimensões Limpas}\\(Bases Geográficas)};

    %% ============================================================
    %% BOUNDING BOXES (Medallion Architecture)
    %% ============================================================
    \begin{scope}[on background layer]
        % Raw / Bronze
        \node[bg_bronze, fit={([shift={(-0.8,0.8)}]api.north west) ([shift={(0.8,-0.8)}]dim_raw.south east)}] (bg_raw) {};
        \node[anchor=north west, font=\sffamily\color{black}, xshift=6pt, yshift=-6pt] at (bg_raw.north west) {\textbf{Camada Bruta / Bronze}};

        % Silver Layer
        \node[bg_silver, fit={([shift={(-0.8,0.8)}]silver_etl.north west) ([shift={(1.0,-0.8)}]silver_data.south east)}] (bg_silver) {};
        \node[anchor=north west, font=\sffamily\color{black}, xshift=6pt, yshift=-6pt] at (bg_silver.north west) {\textbf{Camada Silver / Dados Limpos}};

        % Streaming Processing Pipeline / Gold
        \node[bg_gold, fit={([shift={(-0.8,0.8)}]prod.north west) ([shift={(0.8,0.8)}]topico_focos.north east) ([shift={(0.8,-0.8)}]topico_eventos.south east) ([shift={(-0.8,-0.8)}]buffer.south west)}] (bg_stream) {};
        \node[anchor=north west, font=\sffamily\color{black}, xshift=6pt, yshift=-6pt] at (bg_stream.north west) {\textbf{Pipeline de Streaming / Camada Gold}};
    \end{scope}

    %% ============================================================
    %% SETAS (Roteamento Ortogonal)
    %% ============================================================

    % Raw -> Produtor Kafka
    \draw[relation] (api) -- node[lbl, above=6pt] {streaming} (prod);

    % Produtor -> Fork (bifurcação limpa recuada no eixo X=8.5)
    \draw[relation] (prod.east) -- (8.5, 0) |- (proc_focos.west);
    \draw[relation] (prod.east) -- (8.5, 0) |- (buffer.west);

    % Buffer -> DBSCAN
    \draw[relation] (buffer.east) -- node[lbl, above=6pt] {lote segmentado} (dbscan.west);

    % Cruzamentos com Motor Espacial (tracejado vertical livre)
    \draw[dashed_rel] (proc_focos.south) -- node[right, lbl] {consulta (pontos)} (motor.north);
    \draw[dashed_rel] (dbscan.north) -- node[right, lbl] {consulta (polígonos)} (motor.south);

    % Processamento -> Payloads
    \draw[relation] (proc_focos.east) -- node[lbl, above=6pt] {valida e tipa} (pay_focos.west);
    \draw[relation] (dbscan.east) -- node[lbl, above=6pt] {agrega e tipa} (pay_eventos.west);

    % Payloads -> Tópicos
    \draw[relation] (pay_focos) -- (topico_focos);
    \draw[relation] (pay_eventos) -- (topico_eventos);

    % Raw -> Silver: Fontes Dimensionais descem para ETL
    \draw[relation] (dim_raw.south) |- (silver_etl.west);
    
    % ETL -> Dimensões Limpas
    \draw[relation] (silver_etl) -- (silver_data);

    % Silver -> Motor Espacial (sobe por corredor livre à direita em X=20.0)
    \draw[relation] (silver_data.east) -| node[lbl, near start, right] {carga em memória} (20.0, 0) -- (motor.east);

    \end{tikzpicture}
    }
    \caption{Fluxo de Ingestão e Processamento (Fase 1).}
    \label{fig:diagrama_ingestao}
    \fonte{Elaborada pelos autores.}
\end{figure}


Uma vez inseridos nos tópicos de saída, os dados roteados dão início à segunda fase operacional, cujo consumo analítico é detalhado na Figura \ref{fig:diagrama_analytics}. As mensagens em \textit{streaming} são ingeridas de forma segmentada pelos nós do Apache Druid, consolidando as Tabelas Fato \texttt{fato\_focos} e \texttt{fato\_eventos}. Isso viabiliza pesquisas em frações de segundo e sustenta as renderizações da camada de apresentação, servida através de painéis interativos no Apache Superset.

\begin{figure}[H]
    \centering
    \resizebox{\textwidth}{!}{
    \begin{tikzpicture}[
        font=\sffamily\small,
        %% -- ESTILOS DE COMPONENTES --
        process/.style={rectangle, draw=black!60, thick, fill=white, rounded corners=4pt, drop shadow={opacity=0.15, shadow xshift=2pt, shadow yshift=-2pt}, align=center, inner sep=6pt, minimum width=2.8cm, minimum height=1.3cm},
        database/.style={cylinder, cylinder uses custom fill, cylinder body fill=white, cylinder end fill=gray!10, shape border rotate=90, draw=black!60, thick, drop shadow={opacity=0.15, shadow xshift=2pt, shadow yshift=-2pt}, aspect=0.25, align=center, inner sep=5pt, minimum width=2.8cm, minimum height=1.5cm},
        payload/.style={ellipse, draw=black!60, thick, fill=white, drop shadow={opacity=0.15, shadow xshift=2pt, shadow yshift=-2pt}, align=center, inner sep=4pt, minimum width=2.5cm, minimum height=1.1cm},
        source/.style={tape, tape bend top=none, draw=black!60, thick, fill=white, drop shadow={opacity=0.15, shadow xshift=2pt, shadow yshift=-2pt}, align=center, inner sep=6pt, minimum width=2.8cm, minimum height=1.3cm},
        queue/.style={chamfered rectangle, chamfered rectangle corners={north east, south east}, chamfered rectangle xsep=0.4cm, draw=blue!50, thick, fill=cyan!5, drop shadow={opacity=0.15, shadow xshift=2pt, shadow yshift=-2pt}, align=center, inner sep=6pt, minimum width=2.8cm, minimum height=1.3cm},
        %% -- ESTILOS DE BACKGROUND (MEDALLION) --
        bg_bronze/.style={rectangle, draw=orange!40!brown, thick, dashed, fill=orange!5, rounded corners=6pt, inner sep=12pt},
        bg_silver/.style={rectangle, draw=gray!50, thick, dashed, fill=gray!8, rounded corners=6pt, inner sep=12pt},
        bg_gold/.style={rectangle, draw=yellow!60!orange, thick, dashed, fill=yellow!10, rounded corners=6pt, inner sep=12pt},
        %% -- ESTILOS DE ARESTAS --
        relation/.style={-Stealth, thick, draw=gray!80},
        dashed_rel/.style={-Stealth, thick, dashed, draw=gray!80},
        lbl/.style={font=\sffamily\bfseries\scriptsize, text=black, fill=white, draw=gray!40, rounded corners=2pt, inner sep=3pt}
    ]

    %% ============================================================
    %% GRADE ABSOLUTA (Fluxo Horizontal Left-to-Right)
    %% ============================================================

    % COLUNA 1: INPUT (X=0)
    \node[queue] at (0, 3.5) (topico_focos) {\textbf{Tópico Kafka:}\\Focos};
    \node[database] at (0, 0) (dim_silver) {\textbf{Dimensões Limpas}\\(Camada Silver)};
    \node[queue] at (0, -3.5) (topico_eventos) {\textbf{Tópico Kafka:}\\Eventos};
    
    % COLUNA 2: INGESTÃO (X=6.5)
    \node[process] at (6.5, 3.5) (mid_focos) {\textbf{Nó de Ingestão}\\(MiddleManager)};
    \node[process] at (6.5, 0) (batch_ing) {\textbf{Ingestão Segmentada}\\(Dimensões)};
    \node[process] at (6.5, -3.5) (mid_eventos) {\textbf{Nó de Ingestão}\\(MiddleManager)};
    
    % COLUNA 3: TABELAS ARMAZENADAS (X=13.0)
    \node[database] at (13.0, 3.5) (fato_focos) {\textbf{Tabela Fato:}\\fato\_focos};
    \node[database] at (13.0, 0) (dim_tabela) {\textbf{Tabela Dimensão}\\(Geométrica)};
    \node[database] at (13.0, -3.5) (fato_eventos) {\textbf{Tabela Fato:}\\fato\_eventos};
    
    % COLUNA 4: PRESENTATION / SUPERSET (X=20.0)
    \node[process] at (20.0, 0) (superset) {\textbf{Apache Superset}\\(Painel Visual)};
    
    % COLUNA 5: DASHBOARDS (X=26.0)
    \node[process] at (26.0, 2.5) (dash_risco) {\textbf{Painel}\\de Risco};
    \node[process] at (26.0, -2.5) (dash_terr) {\textbf{Análise}\\Territorial};

    %% ============================================================
    %% BACKGROUNDS (Medallion)
    %% ============================================================
    \begin{scope}[on background layer]
        % Silver Layer Input
        \node[bg_silver, fit={(dim_silver)}] (bg_silver_node) {};

        % Analytical Storage
        \node[bg_gold, fit={([shift={(-0.8,1.5)}]mid_focos.north west) ([shift={(0.8,-0.8)}]fato_eventos.south east)}] (bg_olap) {};
        \node[anchor=north west, font=\sffamily\color{black}, align=left, xshift=6pt, yshift=-4pt] at (bg_olap.north west) {\textbf{Armazenamento Analítico:}\\\textbf{Apache Druid (Camada Gold)}};

        % Presentation Layer
        \node[bg_gold, fit={([shift={(-0.8,1.5)}]superset.west|-dash_risco.north) ([shift={(0.8,-0.8)}]dash_terr.south east)}] (bg_pres) {};
        \node[anchor=north west, font=\sffamily\color{black}, align=left, xshift=6pt, yshift=-4pt] at (bg_pres.north west) {\textbf{Camada de Apresentação}\\\textbf{(Consumo)}};
    \end{scope}

    %% ============================================================
    %% SETAS E ROTEAMENTO ORTOGONAL (.east saindo -> .west entrando)
    %% ============================================================
    \draw[relation] (topico_focos.east) -- (mid_focos.west);
    \draw[relation] (dim_silver.east) -- (batch_ing.west);
    \draw[relation] (topico_eventos.east) -- (mid_eventos.west);
    
    \draw[relation] (mid_focos.east) -- (fato_focos.west);
    \draw[relation] (batch_ing.east) -- (dim_tabela.west);
    \draw[relation] (mid_eventos.east) -- (fato_eventos.west);
    
    % Tabelas -> Superset (Fluxos convergindo ortogonalmente em X=16.5)
    \draw[dashed_rel] (fato_focos.east) -- (16.5, 3.5) |- node[lbl, near start, above=2pt] {SQL} (superset.west);
    \draw[dashed_rel] (dim_tabela.east) -- node[lbl, above=2pt] {SQL} (superset.west);
    \draw[dashed_rel] (fato_eventos.east) -- (16.5, -3.5) |- node[lbl, near start, above=2pt] {SQL} (superset.west);

    % Superset -> Dashboards (Fluxos bifurcando ortogonalmente em X=23.0)
    \draw[relation] (superset.east) -- (23.0, 0) |- (dash_risco.west);
    \draw[relation] (superset.east) -- (23.0, 0) |- (dash_terr.west);

    \end{tikzpicture}
    }
    \caption{Fluxo de Armazenamento Analítico e Visualização (Fase 2).}
    \label{fig:diagrama_analytics}
    \fonte{Elaborada pelos autores.}
\end{figure}


\subsection{Estratégia de Processamento Espacial Híbrido (Python e Druid)}

Embora o Apache Druid possua capacidades nativas de indexação e filtragem espacial, a execução de junções espaciais densas (\textit{Point-in-Polygon joins}) entre milhões de anomalias térmicas e geometrias poligonais complexas, como os limites irregulares de Terras Indígenas ou municípios continentais, representa uma carga computacional intensa. Delegar essa operação de interseção ao motor OLAP durante a ingestão ou consulta degradaria significativamente a vazão (\textit{throughput}) e a latência analítica.

Para mitigar esse gargalo, a arquitetura adota um processamento espacial híbrido, deslocando o enriquecimento geográfico para as camadas prévias de orquestração. O cruzamento é executado em memória através da biblioteca GeoPandas (Python), que tira proveito de estruturas de indexação em árvore (\textit{R-Tree}). Esse método particiona o espaço topológico, reduzindo a complexidade temporal da busca espacial de $O(N)$ para $O(\log N)$. Como resultado, os pacotes (\textit{payloads}) transmitidos via Kafka chegam ao Apache Druid em formato tabular desnormalizado e já enriquecidos com suas chaves infraestruturais, permitindo que o banco atue exclusivamente na agregação analítica de hipervelocidade.
